\section*{Chapitre \ref{ch:s1:crypto-medium-frequency-strategies} --- Stratégies crypto à fréquence moyenne}

\begin{solution}{exo:s1:crypto-medium-frequency-strategies:1}
$0.0001 \times 3 \times 365 = 10.95~\%$ par an~; à 0.1~\% par intervalle, 109.5~\%.
\end{solution}

\begin{solution}{exo:s1:crypto-medium-frequency-strategies:2}
$(30.61~\% - 0.2~\%)/1.2 = 25.3~\%$ sur le capital.
\end{solution}

\begin{solution}{exo:s1:crypto-medium-frequency-strategies:3}
Un mouvement de prix modifie la valeur de la position au comptant et celle du contrat perpétuel vendu, de montants opposés. Le livre est exposé au financement (qui peut devenir négatif),
à la base entre le contrat perpétuel et le comptant, aux appels de marge sur la position vendue lorsque le prix monte (le gain se trouve dans la position au comptant, éventuellement sur une autre plateforme), et
aux plateformes et aux dépositaires qui détiennent les deux jambes.
\end{solution}

\begin{solution}{exo:s1:crypto-medium-frequency-strategies:4}
$0.05 \times 0.808 = 4.0~\%$ par an. Au moins une défaillance~: $1 - 0.95^4 = 18.5~\%$ parmi quatre plateformes, $1 - 0.95^8 = 33.7~\%$ parmi huit (les
18.7~\% et 34.0~\% de la simulation sont du bruit d'échantillonnage autour de ces valeurs).
\end{solution}

\begin{solution}{exo:s1:crypto-medium-frequency-strategies:5}
$25.3\% - 4.0\% = 21.3\%$ en 2021~; $3.3\% - 4.0\% = -0.7\%$ en 2022~: une perte une fois le \omterm{def:s1:crypto-medium-frequency-strategies:venue}{risque de plateforme} intégré.
\end{solution}

\begin{solution}{exo:s1:crypto-medium-frequency-strategies:6}
Une semaine de rendements à 80~\% de volatilité spécifique mesure la dérive injectée avec beaucoup plus de bruit que quatre semaines, si bien que son classement change davantage chaque semaine~:
le livre négocie davantage pour un signal plus faible, et à 20 points de base par unité négociée, les coûts absorbent tout.
\end{solution}

\begin{solution}{exo:s1:crypto-medium-frequency-strategies:7}
Avec la contagion, la perte espérée passe de 4.1~\% à 7.4~\% par an et le 99e centile de 40.4~\% à 60.6~\% du capital~; la probabilité d'une perte sur une
année reste d'environ 18.7~\%, car la contagion n'ajoute de défaillances que les années qui en comptent déjà une. Répartir le capital entre plateformes aide le moins quand les
défaillances surviennent ensemble.
\end{solution}

\begin{solution}{exo:s1:crypto-medium-frequency-strategies:8}
Cinq années sans défaillance de plateforme n'apprennent rien à ce sujet~: le principal risque du livre est de perdre l'essentiel des actifs d'une plateforme, perte que l'historique quotidien des P\&L
ne montre pas avant qu'elle ne survienne. Imputer la perte espérée liée aux plateformes, soumettre le livre à une défaillance simulée, et compter les années de financement faibles ou négatives.
\end{solution}

\begin{solution}{pb:s1:crypto-medium-frequency-strategies:1}
\begin{enumerate}
\item Position longue au comptant et vente d'un contrat perpétuel, pour encaisser le financement~; position longue au comptant et vente d'un contrat à échéance, pour encaisser la base.
\item La prime moyenne du contrat perpétuel par rapport à son indice, à laquelle s'ajoute un terme d'intérêt écrêté, versé toutes les huit heures à 00:00, 08:00 et 16:00 UTC.
\item 14.2~\%, 25.3~\%, 3.3~\%, 6.4~\%, 9.8~\%, 4.1~\% et 2.3~\% pour 2020 à 2026.
\item Il peut dépasser 40~\% par an, varie fortement et reflète la demande d'effet de levier d'investisseurs qui suivent la tendance, face à un capital d'arbitrage limité.
\item Les rendements des crypto-actifs se chargent sur des facteurs propres aux crypto-actifs, avec un fort \omterm{def:s1:trend-following:trend}{momentum temporel}~; le marché, la taille et le momentum captent la dimension transversale.
\item 50 cryptomonnaies, un facteur de marché à 60~\%, une volatilité spécifique de 80~\% à queues épaisses, une dérive injectée de demi-vie de 60 jours.
\item 0.85, 1.56 et 1.29 avant coûts à 7, 28 et 56 jours~; 0.42, 1.35 et 1.13 après 10 points de base~; $-0.01$, 1.13 et 0.97 après 20.
\item Un signal issu des données publiques de la blockchain~; une corrélation injectée de $-0.03$ retrouvée sous forme d'IC de $-0.031$~; et non que des flux réels prédisent les rendements.
\item Le risque de perdre les actifs détenus sur une plateforme ou chez un dépositaire.
\item Des actifs localisés représentant environ un cinquième des dettes envers les clients (un déficit de 80.8~\%), et 7.0~milliards de dollars de retraits nets de clients en onze jours.
\item $80.8~\%/k$ du capital~; 5.0~\%, 9.9~\%, 18.7~\% et 34.0~\% pour une à huit plateformes.
\item La perte espérée vaut $pL$ par unité de capital, où qu'il se trouve.
\item \textbf{Résultat nommé.} Le livre de \omterm{def:s1:crypto-medium-frequency-strategies:carry}{carry de financement} a rapporté 25.3~\% sur le capital en 2021 et 3.3~\% en 2022 (14.2~\% en 2020, 2.3~\% en 2026 jusqu'en septembre)~; réparti
sur quatre plateformes, il perd 20.2~\% du capital lorsqu'une plateforme fait défaut, avec une perte espérée liée aux plateformes de 4.1~\% par an.
\item Choisir les plateformes et la conservation des actifs, garder les actifs hors des bourses et retirer les fonds tôt.
\item Le financement calculé pour chaque intervalle de financement, les frais, le capital réellement immobilisé, les appels de marge et une charge pour perte liée aux plateformes.
\item \omterm{def:s1:crypto-medium-frequency-strategies:carry}{Carry de financement} sur contrat perpétuel détenu sur une seule plateforme.
\item Les deux rapportent le prix de la détention d'un actif que d'autres paient pour s'en délester ou pour y recourir avec effet de levier~; ici, le payeur est la demande à effet de levier.
\item La demande longue à effet de levier sur un marché haussier.
\item La demande à effet de levier a reculé après les pertes de 2022~; le chapitre n'en avance pas la cause.
\item Pour avoir fourni l'effet de levier et supporté le \omterm{def:s1:crypto-medium-frequency-strategies:venue}{risque de plateforme}.
\end{enumerate}
\end{solution}

\begin{solution}{iq:s1:crypto-medium-frequency-strategies:1}
Paiement périodique entre les positions longues et courtes d'un contrat perpétuel, qui n'a pas d'échéance et qui est fixé à partir de la prime du contrat sur son indice~; il ramène le
prix du contrat vers l'indice au comptant en faisant payer le côté le plus chargé.
\end{solution}

\begin{solution}{iq:s1:crypto-medium-frequency-strategies:2}
Acheter au comptant et vendre le contrat perpétuel pour un notionnel égal~; encaisser 60~\% par an tant que cela dure~; les risques sont la baisse ou le passage en négatif du financement, les appels de marge sur la
position vendue si le prix monte, la liquidation si la marge est épuisée, le mouvement de la base et la défaillance de la plateforme~; dimensionner selon le \omterm{def:s1:crypto-medium-frequency-strategies:venue}{risque de plateforme} et non selon le financement.
\end{solution}

\begin{solution}{iq:s1:crypto-medium-frequency-strategies:3}
À partir de la perte que la société peut supporter si une plateforme fait purement et simplement défaut~: un cap par plateforme, modulé selon sa qualité estimée (conservation, audits, juridiction, historique), un
cap sur l'exposition totale aux plateformes sans conservation séparée, et des seuils de retrait (délais de retrait, élargissement des fourchettes, rumeurs).
\end{solution}

\begin{solution}{iq:s1:crypto-medium-frequency-strategies:4}
Positions et marges sur chaque plateforme, financement reçu et versé, soldes et état des retraits, délais de transfert entre plateformes, prix et base
écarts entre eux, et alertes lorsque la marge, le financement ou le comportement d'une plateforme sortent des limites.
\end{solution}

\begin{solution}{iq:s1:crypto-medium-frequency-strategies:5}
Construire l'univers tel qu'il était chaque semaine, y compris les cryptomonnaies radiées par la suite ou éteintes, avec leurs derniers prix négociables~; ne classer que les cryptomonnaies négociables à cette date~;
imputer les coûts cryptomonnaie par cryptomonnaie.
\end{solution}

\begin{solution}{iq:s1:crypto-medium-frequency-strategies:6}
Le nombre de défaillances $N$ suit une loi binomiale de paramètres $(k, p)$ et la perte vaut $NL/k$~: moyenne $pL$, variance $L^2p(1 - p)/k$. La probabilité que la perte soit au moins
$mL/k$ est $P(N \ge m) = \sum_{j \ge m}\binom{k}{j}p^j(1 - p)^{k - j}$.
\end{solution}
